\Needspace{12\baselineskip}
\originalchapter{光电效应与干涉测量}
\originalsection{光电效应的装置和模型}
Junior Lab光电效应实验用汞灯、窄带滤光片、光电管、可调反向电压和灵敏电流计，把入射光频率与电子最大动能联系起来。滤光片选取谱线，透镜将光聚到阴极而尽量避开阳极环；灯的预热会影响谱线强度，暗电流测量必须在相同电子学配置下进行。

一个光子提供能量 $h\nu$。若逸出能为 $\phi$，最大动能 $K_{\max}=h\nu-\phi$，反向截止电压满足 $eV_s=h\nu-\phi$。这条关系预测斜率 $h/e$，不是电流与光强的斜率。增大光强提高发射率，理想条件下不改变最大动能；波长变化才移动截止位置。

实际阴阳极接触电势、杂散光和电流计零点可以改变截距。截止点不是随意挑一个“电流刚好为零”的样本：应对暗电流校正后的一组 $I(V)$ 曲线使用一致外推方法，并检验对取点范围的敏感性。滤光片漏过更短波长的微弱光会人为提高截止电压，属于不易被重复测量发现的系统误差。

拟合 $V_s=A\nu+B$ 后，$h=eA$。若频率来自 $\nu=c/\lambda$，波长不确定度传为 $u_\nu=c u_\lambda/\lambda^2$；同一光谱仪的校准偏差是各点共享项。不能将拟合截距直接当成精确的纯阴极功函数，而忽略接触电势和表面状态。

\begin{center}\begin{tikzpicture}[>=Stealth]
\node[draw] (lamp) at (0,0){汞灯};\node[draw] (filter) at (2,0){滤光片};\node[draw] (tube) at (4.5,0){光电管};
\draw[->] (lamp)--(filter);\draw[->] (filter)--(tube);
\node[draw] (volt) at (3.5,-1.1){反向电压};\node[draw] (current) at (6,-1.1){电流计};
\draw (tube)--(volt);\draw (tube)--(current);
\end{tikzpicture}\end{center}

\originalsection{Michelson干涉仪}
分束器把同一光场分为两臂，各自经镜面反射再叠加。探测强度是 $I=I_1+I_2+2\sqrt{I_1I_2}\cos\Delta\varphi$，相位差含往返路程 $\Delta\varphi=4\pi(L_1-L_2)/\lambda+\varphi_0$。镜面平移 $\Delta L$ 产生 $N=2\Delta L/\lambda$ 个完整条纹，所以 $\lambda=2\Delta L/N$。只算单程会把波长误差整整扩大一倍。

先使两束光斑重合、波前接近同向，再选择能稳定计数的位置。转动镜面改变的是条纹空间分布；平移镜面改变的是整体相位。可见度 $\mathcal V=(I_{\max}-I_{\min})/(I_{\max}+I_{\min})$ 在两束不等强时即使完全相干也小于1。低可见度还可能来自偏振失配、振动或有限相干长度，不能唯一归因于光源谱宽。

误差分析应同时包括微分头标定、机械回差、计数端点和条纹判断。往返扫描、同方向读取、增大总扫描距离可识别不同误差来源；增加条纹数会降低端点误差的相对贡献，却不会降低微分头比例校准误差。

\begin{example}
（编者题）模拟微分头位移 $\Delta L=(.1500\pm.0005)\,\mathrm{mm}$，条纹数 $N=474\pm1$，误差独立。求波长和标准不确定度。
\end{example}
\begin{solution}
$\lambda=2\Delta L/N=632.9\,\mathrm{nm}$，相对方差 $(u_\lambda/\lambda)^2=(u_L/\Delta L)^2+(u_N/N)^2$，得 $u_\lambda\simeq2.5\,\mathrm{nm}$。这个输入是明确的教学模拟读数；若 $\Delta L$ 与计数端点相关，还需协方差项。
\end{solution}

\originalsection{原手册的PZT、光电二极管与示波器读出}
原Optical Interferometry手册的镜位移由压电换能器PZT驱动，微分头读数例题只用于说明几何关系，不能替代这套原测量链。函数发生器同时接PZT及示波器CH2，光电二极管的放大输出接CH1；CH2提供驱动电压 $V_{\rm PZT}(t)$，CH1提供强度的电压表示 $V_{\rm PD}(t)$。对齐两束光在二极管上的位置、使两臂长度接近并选择可用对比度，才能把干涉相位读出来。

在选定扫描支路内采用独立标定 $x(V)=x_0+\kappa_V V$。手册举例器件标签为 $\kappa_V=(41.6\pm1.3)\,\mathrm{nm/V}$，这是手册的标定示例，不是本项目真实器件的测量结果。往返光程给
\[
 V_{\rm PD}=V_{\rm dark}+A[1+\mathcal V\cos(\phi_0+4\pi\kappa_V V_{\rm PZT}/\lambda)].
\]
函数发生器推荐起始设置为10 Hz、10 V峰峰值的正弦或三角波。应记录实际终端电压、示波器探头比例及负载条件，而不是只读发生器面板；信号的同一单调支路可避免把往返扫描中的同相点混为邻近条纹。PZT滞回、蠕变与非线性需要比较上升、下降支路，不能靠增加条纹数消去标定误差。

相邻同相极值的驱动电压间隔为 $\Delta V_{\rm fringe}$，故 $\lambda=2\kappa_V\Delta V_{\rm fringe}$；若用相邻最大与最小，则相位差为 $\pi$，应为 $\lambda=4\kappa_V\Delta V_{\rm half}$。跨N个完整条纹用 $\lambda=2\kappa_V(V_{m+N}-V_m)/N$。也可对一组同相极值拟合 $V_m=V_*+m\Delta V_{\rm fringe}$，检查残差曲率。\emph{驱动电压的极值并不是光电电压的极值}，这两个通道必须按同一时刻配对。

独立校准与间隔读数时，$u_\lambda^2=(2\Delta V_{\rm fringe})^2u_{\kappa_V}^2+(2\kappa_V)^2u_{\Delta V}^2$；共享探头增益或拟合相关性应补协方差项。取手册示例标签，并\emph{另设教学模拟}间隔 $(7.60\pm.03)$ V，实际计算给 $\lambda=632.32$ nm、$u_\lambda=19.92$ nm，主要误差来自PZT标签的独立标定。没有将此模拟读数称为原实验观测。

手册分析题要求求可见度并解释非理想性、由电压校准推断波长。先减去暗电压后，$\mathcal V=(V_{\max}-V_{\min})/(V_{\max}+V_{\min}-2V_{\rm dark})$；两臂强度不等、偏振不一致、波前失配和相干长度限制都可降低它。记录极值时间、电压、支路方向和器件标签，保存原始双通道轨迹，才能分别追溯随机读数与共享校准误差。

\begin{center}\begin{tikzpicture}[>=Stealth,node distance=1cm]
\node[draw](g) at (0,0){函数发生器};\node[draw](p) at (3,0){PZT与镜};\node[draw](d) at (6,0){光电二极管};
\draw[->](g)--node[above]{驱动电压}(p);\draw[->,dashed](p)--node[above]{干涉强度}(d);
\node[draw](c2) at (0,-1.1){示波器CH2};\node[draw](c1) at (6,-1.1){示波器CH1};\draw[->](g)--(c2);\draw[->](d)--(c1);
\end{tikzpicture}\end{center}

\originalsection{光电截止点的推断及误差预算}
发射电子通常具有动能分布而非都在最大能量处。理想一维示意中，反向势 $V$ 截走 $K<eV$ 的电子，$I(V)\propto\int_{eV}^{K_{\max}}f(K)\dd K$；真实光电管还包含角度接受与阴阳极几何。因此取电流零交点、尾部线性外推或某固定分数电流的方法会给不同系统偏移。手册要求保留实际曲线，是因为这个推断环节不能由爱因斯坦直线公式代替。

可按以下连续流程采集和分析：稳定光源与电子学；遮光测同量程暗曲线；用一条谱线从饱和电流区向截止区扫描；在截止附近增加读数并反向重复以检查滞后；记录光斑是否照到阳极；换滤光条件重复相同规则。采集顺序可交错波长，以免灯漂移与频率完全共线。此处描述功能流程，没有虚构某次实验已经完成。

对于 $V_s=A\nu+B$，频率零点和比例标定分别影响截距和斜率，电压零点主要影响截距、电压比例影响斜率。接触势若对波长不变，也主要进截距；漏光可能改变曲线形状和截止提取，不一定能用一个共同常数吸收。故“斜率法可消去截距”不能解释为全部系统误差都消失。

附带教学模拟取六个频率、独立截止标准差 $.012$ V、固定随机种子20261012。实际脚本得 $h_{\rm fit}=6.54115\times10^{-34}\,\mathrm{J\,s}$、标准不确定度 $9.98395\times10^{-36}\,\mathrm{J\,s}$，$\chi^2=1.7564$、自由度4。这只是合成误差模型下的一次实现，和SI定义常数比较在一个标准差内；没有因此声称Junior Lab仪器已测得该值。模拟文件允许改变共享零点与比例误差，观察二者如何不同地影响拟合。

\originalsection{干涉仪的光程和相干条件}
光场沿一臂经过折射率 $n(s)$ 的路径，相位为 $2\pi\int n(s)\dd s/\lambda_0$。Michelson有往返两程，若长度 $\ell$ 的气室折射率由 $n_1$ 改到 $n_2$，固定镜面时条纹数 $N=2\ell(n_2-n_1)/\lambda_0$。镜面运动测几何位移，气室变化测光学程，二者都来自同一个相位差。

有两条接近谱线 $\lambda_1,\lambda_2$ 时，强度干涉项相加形成包络。用波数 $k_i=2\pi/\lambda_i$，两项和为 $2\cos[(k_1-k_2)\Delta L]\cos[(k_1+k_2)\Delta L]$（此处 $\Delta L$ 为单臂差）。相邻可见度极小间的单臂差约 $\Delta L_{\rm beat}=\lambda^2/(2|\Delta\lambda|)$。先明确横轴单程还是总光程，才能决定因子二。

宽谱光源的可见度随时延衰减，等于归一化场相关函数的模。频谱与相关函数互为Fourier变换，窄谱通常相干时间长；但偏振或强度不匹配也降低可见度。对等强、同偏振条件下的Gaussian频谱，相关包络仍为Gaussian，这与单一指数寿命展宽给出的指数包络不同。

\originalsection{机械回差与条纹计数的具体判断}
微分头的读数差不是镜子位移的无条件真值：传动比、杠杆几何和弹性变形可能进入标定。每段扫描沿同一方向逼近读数，往返全程用来检查方向相关偏差。若端点条纹判断约 $\pm1/2$ 周而两端独立，差值计数标准不确定度是相应合成值；用 $\pm1$ 周只是明确的保守计数约定，并非由数据自然得出。

平移次数很多时可将累计条纹数对位移拟合，减少两端点依赖；然而累计计数共享前面所有误判，残差会相关，不能假装各累计点独立。也可把若干不重叠位移段分别记录。无论采用哪一种，图和文件必须能对应到所选方法。
